\begin{minipage}[b]{2.5in}
  Resubmission Cover Letter \\
  {\textbf GENETICS}
\end{minipage}
\hfill
\begin{minipage}[b]{2.5in}
    Jonathan Terhorst  \\
  \today
\end{minipage}

\vskip 2em

\noindent
{\bf To the Editor(s) -- }

\vskip 1em

We are writing to submit a revised version of our manuscript,
titled
``Parameterizing the genetic architecture under stabilizing selection'',
GENETICS-2026-309276.


\vskip 2em

\noindent \hspace{4em}
\begin{minipage}{3in}
\noindent
{\bf Sincerely,}

\vskip 2em

{\bf
   Jonathan Terhorst, on behalf of the authors.
}\\
\end{minipage}

\vskip 4em

\pagebreak

%%%%%%%%%%%%%%
\reviewersection{1}

% Major comments:

\begin{point}{} % \revpoint{1}{1} in paper.tex
The authors actually do two distinct things – they present a general method to go from a
population genetic model to a mixed effect model. That is, given a distribution of
selection effect $p(s)$, a distribution of allele frequencies $p(x\mid s)$, and the expectation of
the squared effect size given the selection coefficient $\E{\beta^2 \mid s}$ they get
\begin{align*}
    \E{\beta^2 \mid x} = 
    \frac{\int \E{\beta^2 \mid s}p(x \mid s) p(s)ds}{\int p(x \mid s) p(s) ds}
\end{align*}
While this seems like a trivial and unexciting result, it is vital for formulating population
genetic models (such as the Simons model) as mixed effects models. Therefore, seeing
it explicitly explained and justified in print is extremely important.
The second thing done by the author is present a specific model to which they apply
this method. They assume that $p(s)$ is a chi-squared distribution, $p(x \mid s)$ follows its
constant-population-size under-dominant-selection low-mutation-rate form, and suggest
a new form for $\E{\beta^2 \mid s}$ as a mix of constant effect and selection dependent effect
(similar to Simons et al). These assumptions, together, result in equation 19.
I think the paper needs a restructuring in which the general method is separated from
the particular model. This will greatly strengthen the paper because it will emphasize
that future work can make other assumptions and derive analytic (or numeric) versions
of equation 19 for their model. The simulation and inference section can then serve as a
proof of concept of how to go from pop-gen model to mixed model to inference.
\end{point}

\reply{
placeholder
}

\begin{point}{} % \revpoint{1}{2} in paper.tex
There is an obvious gap in developing such a model, running simulations to show its
use, but not applying it to data. One option is to apply it to UK biobank data and
compare the results to alpha model results. However, that would very much change the
focus of the paper. The alternative is to use the discussion to explain why they did not
apply it to data and what do they expect to see if it were fit to data.
\end{point}

\reply{
We did not conduct a full analysis on human genomes because an efficient implementation of the proposed model warrants a separate methods paper.
}

% Minor comments:

\begin{point}{} % \revpoint{1}{3} in paper.tex
The introduction starts quite abruptly. I feel there is a need for an introductory
paragraph.
\end{point}

\reply{
Placeholder
}

\begin{point}{} % Line 152, \revpoint{1}{4} in paper.tex
Notation. The authors introduce much notation and maybe some simplification could be
achieved. Some examples:
\begin{enumerate}
    \item[(a)] Allele frequency is represented as both x and p. This is particularly confusing in
the integral in section 2.4 where both are used in the same equation and p is also
used to denote probability distributions.
    \item[(b)] Beta and alpha are both effect sizes and it’s not completely clear what the
distinction is.
    \item[(c)] Is $\alpha$ actually needed? Can it not be replaced by $s_j = \frac{\norm{\alpha_j}^2}{W_S}$?
\end{enumerate}
\end{point}

\reply{
\begin{enumerate}
    \item[(a)] We have replaced all $p$ with $x$ to denote allele frequencies.
    \item[(b)] We added a paragraph clarifying their definitions and mutual relationship.
    \item[(c)] The selection coefficient is a function of the allele frequency in stabilizing selection.
    We also wanted to emphasize that the latent trait $z$ is a quantitative trait as a function of $\alpha_j$.
\end{enumerate}
}

\begin{point}{} % \revpoint{1}{5} in paper.tex
$\E{\beta^2 \mid s}$ (equation 13) kind of drops out of thin air. It would be better to justify it before it
appears. One could first introduce the simple model of equation 16, and then suggest to
use the resulting $\E{\beta^2 \mid s}$ more generally. One could also motivate it by a probabilistic
model in which each variant has either no pleiotropic fitness effects with probability $\rho^2$ or
has pleiotropic fitness effects (and behaves similarly to Simons et al with $\beta^2 \propto \, s$ with probability $\rho^2$.
\end{point}

\reply{
placeholder
}

\begin{point}{}% \revpoint{1}{6} in paper.tex
When doing that it would be a great opportunity to give some intuition into the meaning
of $\sigma_a^2$, $\sigma_b^2$, and $\rho^2$.
within section 2.3. Such intuition does appear downstream, but, on first
read, I felt it was sorely needed in section 2.3.
\end{point}

\reply{
placeholder
}

\begin{point}{} % \revpoint{1}{7} in paper.tex
In equation 13, I think it’s better to write the right hand term as $\rho^2 \cdot \sigma_b^2 \cdot \frac{\norm{\alpha}^2}{\sigma_a^2}$.
Since $\E{\frac{\norm{\alpha}^2}{\sigma_a^2}}=1$
this emphasize that $\sigma_b^2$
determines the scale of effect sizes and the coupling
to selection just effect how they depend on selection.
\end{point}

\reply{
   Placeholder. 
}

\begin{point}{} % Line 313, \revpoint{1}{8} in paper.tex
Can we get a figure that shows the behavior of equation 19? I had to plot it myself to get
a sense of it.
\end{point}

\reply{
Figures~\ref{fig:blup_beta2_by_maf_alpha_1},~\ref{fig:blup_beta2_by_maf_alpha_0p5_appendix} and~\ref{fig:blup_beta2_by_maf_alpha_0_appendix} draw the requested curve alongside two $\alpha$-model variants with or without MAF-binning.
}

\begin{point}{} % Figure 4, % \revpoint{1}{9} in paper.tex
Speaking of which, the plots I plotted look very similar to Supplementary Figure 4 of
Schoech 2019 (which is often overlooked and should have been the basis of an
independent paper). Since this figure was trying to show that the alpha model cannot
describe realistic genetic architectures, it would be nice to show that this model can
achieve that and can match the architectures simulated by Schoech.
\end{point}

\reply{
This would be a great addition.
Figures~\ref{fig:blup_beta2_by_maf_alpha_1},~\ref{fig:blup_beta2_by_maf_alpha_0p5_appendix} and~\ref{fig:blup_beta2_by_maf_alpha_0_appendix} are precisely the analogous plot to Schoech's supplementary figure.
}

\newpage

%%%%%%%%%%%%%%
\reviewersection{2}

\begin{point}{} % \revpoint{2}{1} in paper.tex
I find the results presented in figures 3 and 4 somewhat confusing. Both indicate that the alpha=0 model implying no relationship between effect size and frequency is closer to the simulated model of selection than either alpha>0 model. This demands a clear explanation given that when alpha models have been fit in practice estimates are consistently significantly greater than zero. Moreover, the simulation results also show particularly poor performance for alpha=1. Given that this is the standard parameterization in the field, and methods do not perform exceedingly poorly, it begs the question of whether the simulations are particularly realistic, at least for humans. The explanation in the discussion that the good performance of alpha=0 is partly due to rare variants where the dependency disappears is also concerning because (again, in humans) rare variants don't contribute massively to heritability. I wonder if the restriction to simulating a single trait under stabilizing selections is contributing to some of this. This all needs a clearer explanation before a strong preference for evolutionary derived models can be claimed off of the results.
\end{point}

\reply{
The reviewer raised an important question of reconciling the manuscript's findings against the $\alpha$-model and the model's empirical success.
We argue that minor allele frequency (MAF) stratification commonly shipped together with the $\alpha$-model improves the model fit at the cost of increased computational burden, lower precision, and lost interpretability of parameters.
This approach is standard with examples including stratified LDSC (Gazal et al. 2017, Nature Genetics), GREML-LDMS (Yang et al. 2015, Nature Genetics), and LDAK (Speed et al. 2020, Nature Genetics).
These models adopt a common $\alpha$ value across MAF bins but separately estimate the numerator $\sigma_b^2$ for each bin.
Therefore, they consist of an ensemble of $\E{\beta^2 \mid p} =\sigma_b^2 \,/\, [p(1-p)]^{\alpha}$ curves (as a function of MAF $p$) instead of a single $\alpha$-model representing the whole frequency spectrum.

Figure X, X, and X plot three different $\E{\beta^2 \mid p}$ parameterizations against MAF $p$: our evolutionary model, the vanilla $\alpha$-model, and the MAF-binned $\alpha$-model.
The MAF bins in the figure was adopted from Yang et al. (2015, Nature Genetics).
The true $\E{\beta^2 \mid p}$ curve represented by the evolutionary model drops gradually with increasing MAF.
The vanilla $\alpha$-model deviates significantly from the evolutionary model due to $[p(1-p)]^{-\alpha}$ diverging to infinity for rare variants.
The MAF-binned $\alpha$-model suffers less from the same problem because all MAF bins except for the lowest MAF have bounded values.
Furthermore, within MAF bins, the curves closely resemble the true evolutionary curve.
As shown in Figure X, X, and X, this translates to equivalent prediction accuracy to that of the evolutionary model.

This is not at all surprising.
By finely dividing the MAF spectrum and fitting separate $\E{\beta^2 \mid p}$ per MAF bin,
we can approximate any continuous and bounded curves of the true underlying $\E{\beta^2 \mid p}$ function.
This can be done with step functions (corresponding to $\alpha=0$) or $\alpha$-models of different positive $\alpha$ values ($\alpha=0.5, 1$) as shown in Figure X, X, and X.
In fact, as long as the MAF bins are fine enough, one may choose any parametric curve family -- e.g. linear, quadractic, or exponential -- for each MAF bin because any two continuous curves are indistinguishable for a short enough interval.
That is, the common practice of MAF stratification is a highly flexible model of frequency-dependent architecture that is insensitive to the choice of $\alpha$: it is the MAF binning doing the heavy lifting.

We also examined the cumulative genic variance -- the numerator of SNP-heritability $h^2$ -- across the MAF spectrum in Figure X-X.
The vanilla $\alpha$-model does very badly, aligning with its poor predictive performance whereas the MAF-binned $\alpha$-model follows the evolutionary curve almost perfectly.
As pointed out by the reviewer, rare variants occupy only a small fraction of the total genic variance, so the poor $\E{\beta^2 \mid p}$ fit in rare MAF bins in Figure X, X, and X, would not have mattered much in prediction tasks.

Therefore, MAF-stratified $\alpha$-model performs nearly as good as the correct evolutionary model in prediction (Figure X).
This explains the empirical success of the $\alpha$-model: 
MAF-stratified $\alpha$-model with sufficiently fine MAF bins behave as a non-parametric curve estimator that can adapt to any frequency-dependent architecture.
As prediction only cares about how well the model fits to data, 
MAF-stratified $\alpha$-model is a good choice for this purpose.
Nevertheless, this advantage comes at the cost of lost interpretability.

The discussion so far addresses the reviewer's questions as follows.
The simulation setting, as shown in Figure X-X, reflects the low contribution of rare variants to heritability.
It also demonstrates $\alpha$-model's good fit to data with MAF binning as shown in earlier human genetics studies.
Furthermore, the positive $\alpha$-value reported in the literature does not contradict with our finding.
The vanilla $\alpha$-model fits data better with $\alpha=0$ (Figure X-X), 
but the more common MAF-binned $\alpha$-model matches the evolutionary curve better with $\alpha=0.5$ than $\alpha=0$ (Figure X-X).

We have updated Results and Discussion to reflect these results.

}

\begin{point}{}{}
Also, Schraiber et al 2024 showed that alpha=1/standard GRM corresponded to mutation-selection balance, and this should be noted. 
\end{point}

\reply{
Connecting our mutation-selection-drift balance result to the deterministic ($N_e \rightarrow \infty$) mutation-selection balance is a great exercise.
In fact, Schraiber et al. (2024, PLoS Biology)'s result is a special case of our formula with $N_e \rightarrow \infty$ up to a constant factor of mutation rate $\mu$.
To see this, modify equation (19):
\begin{align*}
    \begin{aligned}
        \E{\beta^2 \mid p} &=
        \frac{\sigma_a^2}{1 + 2\frac{\sigma_a^2N_e}{V_S}p(1-p)} \\
        &\approx 
        \frac{V_S}{2N_e p(1-p)} \\
        &\propto \, \frac{1}{p(1-p)}
    \end{aligned}
\end{align*}
with $W_S = V_S / (2N_e)$.


}

\begin{point}{} % Line 658, \revpoint{2}{2} in paper.tex
Previous work on this topic noted the importance of the DFE (e.g. Schoech et al.) in shaping the effect given frequency function. The authors should note explicitly that the prior term in eq 17 is the DFE. Given the demonstrated importance of the DFE in prior work the sensitivity of 19 to the assumed chi2 DFE should be investigated.
\end{point}

\reply{
placeholder
}

\begin{point}{} % \revpoint{2}{3} in paper.tex
It would seriously improve reader understanding if eq 19 were plotted under a reasonable range of scenarios and compared to alpha models.
\end{point}

\reply{
Figure X-X shows the equations alongside two different $\alpha$-models with and without MAF-binning.
}

\begin{point}{} % \revpoint{2}{4} NOT NEEDED in paper.tex
It isn't really clear what eqs 3 and 4 actually add here given that eq 6 can be obtained from drift and selection terms alone.
\end{point}

\reply{
Based on Bulmer's paper provided by the reviewer,
the suggested derivation likely refers to the use of Wright's formula.
Wright's formula (equation 21 of the same Bulmer reference).
Our derivation using Kolmogorov's equation is merely an elaboration of Bulmer as Wright's formula is proved from Kolmogorov's equation.
}

\begin{point}{}
   The derivation of eq 7 parallels Bulmer 1972, which should be cited. 
\end{point}

\reply{
Thank you for pointing to the reference.
We've updated the reference.
}

\begin{point}{}
Given that h2 is intended to be an estimator of the narrow sense heritability explained by SNPs, and that this is a quantity of considerable interested, it should be noted whether such a quantity can also be derived from your model.
\end{point}

\reply{
We explain how genic variance analogous to $h^2$ can be obtained from our model.
The same logic was used to draw Figure X-X.
}


\begin{point}{}
More care needs to be taken to explain the distinction between alpha and beta for the focal and pleiotropic traits.

\end{point}

\reply{
placeholder
}

\begin{point}{}
Eq 13 needs greater explanation. The manuscript states that the key is to deduce the relationship between effect and frequency, but 13 feels ad hoc as a generalization to the Simons et al. model. Models of this nature also have a tendency to produce a genetic architecture with large effect variants at high frequencies.
\end{point}

\reply{
placeholder
}

\begin{point}{}
Does the derivation leading to eq 18 include a large sample size limit? It isn't clear from the text.

\end{point}

\reply{
Appendix
}
